\subsection{固定代数体例の全中心・素性・PI次数9}
\label{subsec:cubic-third-root-unit-boundary-L-finite-CSA}
\begin{theorem}[階数81の基準例の全中心と中心単純次数]
\label{thm:cubic-third-root-unit-boundary-L-finite-CSA}
前節の固定代数体例では
\[
 Z(A)=K[x,\Psi,z_9],\qquad
 A\otimes_{K[x,\Psi,z_9]}K(x,\Psi,z_9)
 \text{ は次数9の中心単純代数}
\]
であり、$A$ は素である。
標数0代数閉体 $k$ への任意の埋込み後も、
アフィン切断 $U$ は素、$Z(U)=k[\psi,z]$、PI次数9となる。
ここで $z=z_9x^{-9}$。
その中心単純代数のBrauer類を $\alpha$ とすると
\[
 R\text{ が整域}\ \Longleftrightarrow\ U\text{ が整域}
 \ \Longleftrightarrow\ \operatorname{per}(\alpha)=9.
\]
周期が1,3,9のどれか、H3許容性、全パラメータ版は未決定である。
\end{theorem}
\begin{proof}
前節より $A$ は $S=K[x,\Psi,z_9]$ 上自由階数81。
中心fiber $x=1,\Psi=0,z_9=1$ を調べる。
このfiberは81次元で、斉次商
\[
 A/(\Psi,\ z_9-x^9)
\]
の次数12及び13は各81次元。
$x$ 倍が同型となるので、次数12基底からfiberの全乗法を回収できる。

係数の有限生成整環を適当に局所化して還元
\[
 109=0,\qquad t\mapsto15,\qquad\omega\mapsto63
\]
を取る。これは標数0の体から有限体への体準同型という意味ではなく、
実際に使用する全係数の分母が109で可逆な整数環の還元である。
最小多項式と $L,\Omega$ の条件を直接確認する。
この還元で、次数12の全81基底に対する交換子行列のrankは80。
さらに正則表現のtrace形式
\[
 T_{ij}=\operatorname{Tr}(R_{a_i a_j})
\]
の $81\times81$ Gram行列式は $45\bmod109$ となり非零である。
元の六二次関係、三次切断、中心fiber二関係を全生成元行列へ再代入して零を確認する。
Gram行列の対称性も検証する。

用いた基底のpivot、rank80のminor、trace行列式は全て非零のwitnessであり、
標数0でも非零である。
一方、fiber次元81という上限は有限体計算から推測せず、前節の標数0の自由加群定理から既知。
従って同じ基底が標数0へ持ち上がり、そのfiberの中心は一次元、trace形式は非退化となる。
trace形式の核はJacobson radicalを含むので、このfiberは半単純である。
同じ非零minorと行列式を中心一般点へ移せば、
$A_F$（$F=K(x,\Psi,z_9)$）も中心一次元・半単純、従って中心単純。
$\dim_F A_F=81$ より次数は9。

全中心の元は $A$ の有限 $S$ 加群性から $S$ 上整であり、
一般点では $F$ に入る。
$S$ は整閉な多項式環なので $Z(A)=S$。
$A$ は $S$ 上torsion-freeで一般点が単純なので、
任意の二つの非零両側idealの積は非零、すなわち $A$ は素。
中央 $x$ の局所化と次数零部分を取ると、$U$ の全中心と素性、PI次数9が従う。
これらの非零rank/中心単純性は体拡大で保たれるので、全六埋込みと $k$ への拡大に有効である。

最後に $k(\psi,z)$ は代数閉体上超越次数2なので、
period--index定理 \cite[Theorem~4.2.2.3]{LieblichPeriodIndex}を適用できる。
次数9の中心単純代数が除法代数であることはindex9に同値、従って周期9に同値。
$U$ は中心一般点に単射するorderなので、その一般点が除法であることと $U$ が整域であることは同値。
$R$ と $U$ の整域性の同値は既証である。

検証器はNumPyを使うが、演算は全てmod109の整数行列であり、浮動小数点は使わない。
各行列積の内積上限 $81\cdot108^2=944784<2^{63}$ を確認してからint64で計算し、
各積の直後に還元する。
これは精度を仮定する数値実験ではなく、明示的な有限体の非零証明書である。
\end{proof}
